% VI. Експериментални резултати и анализ.
% Всяко число тук идва от записаните изпълнения в results/ на хранилището.
% Числа, които не са измерени, не се появяват.
\section{Експериментални резултати и анализ}
\label{sec:results}

Разделът представя резултатите по методиката от Раздел~\ref{sec:method}. Редът следва реда
на въпросите: първо съвпадението с еталона, защото без него измерването на времето няма
смисъл, после цената на заявката, после цената на останалите стъпки.

Всички числа идват от два записа: \code{results/measurements\_baseline.txt} за
подразбиращата се конфигурация и \code{results/measurements\_native.txt} за конфигурацията
с \code{-march=native}. Съответните файлове във формат CSV са до тях.

\subsection{Съвпадение с еталона}

Наборът от проверки съдържа 60 случая, разпределени в шест групи, и всички минават. Част
от тях са изброени в Табл.~\ref{tab:correctness} заедно с еталона, срещу който сравняват.

\begin{table}[H]
\centering
\caption{Проверка за коректност: какво се сравнява и срещу какво}
\label{tab:correctness}
\begin{tabular}{@{}L{3.6cm}L{4.6cm}L{5.4cm}@{}}
\toprule
\textbf{Проверявано} & \textbf{Еталон} & \textbf{Резултат} \\
\midrule
Заявка за $k$ най-близки съседи & изчерпателно търсене, $k = 1, 5, 32$, 40 заявки & съвпадение по множеството от разстояния \\
Скаларен срещу пакетен път & един срещу друг, 60 заявки & съвпадение до бит, допуск 0 \\
Заявка по радиус & изчерпателно търсене, три радиуса, 30 заявки & съвпадение по множеството от номера \\
Заявка за единствен съсед & първият резултат от заявката за $k$ & съвпадение до $10^{-6}$ \\
Размер на листа & дърво с лист 1 срещу лист 4, 16, 64, 256 & еднакъв отговор \\
Изродени облаци & изчерпателно търсене & еднакъв отговор, построяването приключва \\
Сингулярно разлагане & $USV^{T}$ срещу входа & норма на разликата под $10^{-10}$ \\
Оценяване по Кабш & зададено преобразувание & ъгъл под $10^{-5}$ rad, транслация под $10^{-4}$ \\
Кабш върху равнина & детерминанта на ротацията & в интервала $(0{,}99;\ 1{,}01)$, не отражение \\
Съвместяване без шум & зададено преобразувание & ъгъл под $5 \cdot 10^{-3}$ rad \\
Кръговрат PLY, PCD, PGM, PPM & входният облак или кадър & съвпадение до $10^{-5}$, а за кадрите до бит \\
\bottomrule
\end{tabular}
\end{table}

Съвпадението на двата пътя е с допуск точно нула, а не приблизително. Това е по-силно
твърдение от „дават близки резултати“ и е нарочно: двата пътя извършват едни и същи
операции с плаваща запетая в един и същ ред, само групирани различно, затова всяко
разминаване би било грешка, а не закръгляване.

Проверката \code{icp.kabsch\_never\_returns\_a\_reflection} улови действителна грешка в
сингулярното разлагане, описана в Раздел~\ref{sec:implementation}. Тя е показателна за
това защо еталонът трябва да е независим: реализацията изглеждаше вярна и минаваше
проверката за възстановяване на матрицата, а грешеше само при вход с ранг под три.

\subsection{Построяване на индекса}

Табл.~\ref{tab:build} дава времето за построяване и заетата памет. Броят възли и паметта
не зависят от товара на машината и се повтарят точно.

\begin{table}[H]
\centering
\caption{Построяване на k-d дървото, размер на листа 32, 5 повторения}
\label{tab:build}
\begin{tabular}{@{}rrrrrr@{}}
\toprule
\textbf{Точки} & \textbf{Мин., ms} & \textbf{Мед., ms} & \textbf{Разс., \%} & \textbf{Памет, MiB} & \textbf{Възли} \\
\midrule
10\,000    &   1{,}379 &   1{,}477 &  7{,}9 &  0{,}18 &  1\,023 \\
50\,000    &   8{,}440 &   8{,}855 & 11{,}8 &  0{,}88 &  4\,095 \\
200\,000   &  50{,}659 &  52{,}831 & 10{,}6 &  3{,}53 & 16\,383 \\
1\,000\,000 & 432{,}459 & 575{,}300 & 59{,}3 & 17{,}64 & 65\,535 \\
\bottomrule
\end{tabular}
\end{table}

Паметта расте точно линейно: 17,64 MiB за милион точки са 18,5 байта на точка, при 12
байта координати, 4 байта номер и остатъка за възлите. Времето расте по-бързо от
линейното: от 50\,000 до 1\,000\,000 точки, тоест двадесеткратно увеличение, времето
нараства петдесеткратно. Част от разликата е очакваната от $n \log n$ на
\code{nth\_element}, но множителят е по-голям и последният ред е с разсейване 59,3\,\%,
което го прави най-ненадеждния в таблицата. При \code{-march=native} същият случай дава
352,036 ms, тоест конфигурацията на компилиране влияе и на построяването, макар в него да
няма векторизиран цикъл; разликата е в границите на разсейването и не се тълкува.

\subsection{Цена на заявката за съседи}

Табл.~\ref{tab:knn} дава времето за 20\,000 заявки при трите стойности на $k$.

\begin{table}[H]
\centering
\caption{Заявка за $k$ най-близки съседи, 20\,000 заявки, 9 повторения, подразбираща се конфигурация}
\label{tab:knn}
\begin{tabular}{@{}rrrrrrr@{}}
\toprule
\textbf{Точки} & \textbf{$k$} & \textbf{Скал., ms} & \textbf{Пак., ms} & \textbf{Отн.} & \textbf{Разс., \%} & \textbf{Заявки/s} \\
\midrule
10\,000 &  1 &   8{,}70 &   8{,}48 & 1{,}03 &  16{,}2 & 2\,357\,823 \\
10\,000 &  8 &  26{,}81 &  25{,}13 & 1{,}07 &  54{,}7 &    795\,909 \\
10\,000 & 32 &  85{,}13 &  90{,}71 & 0{,}94 & 171{,}2 &    220\,495 \\
50\,000 &  1 &  12{,}20 &  11{,}40 & 1{,}07 & 169{,}1 & 1\,753\,986 \\
50\,000 &  8 & 111{,}40 &  41{,}47 & 2{,}69 & 215{,}4 &    482\,243 \\
50\,000 & 32 & 132{,}98 & 107{,}22 & 1{,}24 &  61{,}9 &    186\,535 \\
200\,000 &  1 &  17{,}76 &  18{,}50 & 0{,}96 & 103{,}4 & 1\,080\,818 \\
200\,000 &  8 & 134{,}36 &  50{,}07 & 2{,}68 & 147{,}9 &    399\,463 \\
200\,000 & 32 & 136{,}71 & 148{,}32 & 0{,}92 & 107{,}9 &    134\,841 \\
1\,000\,000 &  1 &  31{,}44 &  31{,}95 & 0{,}98 &  48{,}8 &    626\,047 \\
1\,000\,000 &  8 &  65{,}09 &  58{,}88 & 1{,}11 &  38{,}5 &    339\,655 \\
1\,000\,000 & 32 & 161{,}22 & 139{,}85 & 1{,}15 &  40{,}6 &    143\,006 \\
\bottomrule
\end{tabular}
\end{table}

Таблицата не позволява извод за отношението между двата пътя и това е самият резултат.
Отношението се движи между 0,92 и 2,69, а разсейването в същите редове стига 215\,\%.
Двата реда с отношение около 2,69, при 50\,000 и 200\,000 точки и $k = 8$, са придружени
от най-голямото разсейване в таблицата: там минимумът на скаларния път не е уловил
неразтревожено изпълнение. Тези редове не се тълкуват. Същите случаи в конфигурацията с
\code{-march=native} дават 0,88 и 1,13, което потвърждава, че стойността около 2,69 е
следствие от товара на машината, а не от кода.

Изводът, който данните поддържат, е отрицателен и е записан като такъв: \term{разликата
между двата пътя в цялата заявка е под разделителната способност на тази машина}. Част от
времето на заявката отива в обхождане на дървото и в работа с купчината, а те са
еднакви и за двата пътя. Колко точно е тази част, се вижда от следващото измерване.

Пропускателната способност е величина, за която таблицата все пак стига: около 2,4
милиона заявки в секунда при $k = 1$ върху 10\,000 точки и около 143\,000 при $k = 32$
върху милион точки. Спадът от 10\,000 до 1\,000\,000 точки при $k = 32$ е под два пъти,
при стократно увеличение на облака, което съответства на логаритмичната зависимост,
очаквана от устройството на дървото.

\subsection{Изолирано обхождане на лист}

Табл.~\ref{tab:leafscan} премахва обхождането на дървото от измерването: дървото е
построено с размер на листа над броя на точките, тоест е един лист, и всяка заявка е само
обхождане на лист. Общата работа е изравнена между редовете.

\begin{table}[H]
\centering
\caption{Изолирано обхождане на лист, 9 повторения. Ляво: подразбираща се конфигурация. Дясно: с \code{-march=native}}
\label{tab:leafscan}
\begin{tabular}{@{}rrrrrrrrr@{}}
\toprule
& & \multicolumn{3}{c}{\textbf{Подразбираща се}} & & \multicolumn{3}{c}{\textbf{\code{-march=native}}} \\
\cmidrule(lr){3-5}\cmidrule(lr){7-9}
\textbf{Точки в листа} & \textbf{$k$} & \textbf{Скал., ms} & \textbf{Пак., ms} & \textbf{Отн.} & &
\textbf{Скал., ms} & \textbf{Пак., ms} & \textbf{Отн.} \\
\midrule
1\,024  & 1 &  93{,}76 & 37{,}18 & 2{,}52 & &  54{,}12 & 23{,}77 & 2{,}28 \\
1\,024  & 8 & 109{,}85 & 82{,}44 & 1{,}33 & &  91{,}70 & 70{,}46 & 1{,}30 \\
8\,192  & 1 &  55{,}79 & 26{,}16 & 2{,}13 & &  53{,}95 & 17{,}72 & 3{,}05 \\
8\,192  & 8 &  62{,}63 & 30{,}99 & 2{,}02 & &  57{,}23 & 24{,}68 & 2{,}32 \\
65\,536 & 1 &  52{,}91 & 22{,}23 & 2{,}38 & &  49{,}87 & 15{,}30 & 3{,}26 \\
65\,536 & 8 &  42{,}97 & 18{,}38 & 2{,}34 & &  55{,}25 & 23{,}24 & 2{,}38 \\
\bottomrule
\end{tabular}
\end{table}

Тук резултатът е еднозначен. Пакетният път е между 2,0 и 2,5 пъти по-бърз в
подразбиращата се конфигурация и между 2,3 и 3,3 пъти при \code{-march=native}, ако се
изключи редът с $k = 8$ и лист от 1\,024 точки. Посоката е една и съща във всичките
дванадесет клетки, при две независими конфигурации на компилиране.

Печалбата от по-широкия регистър е видима и подредена: при лист от 8\,192 и 65\,536 точки
отношението нараства от 2,13 и 2,38 на 3,05 и 3,26 при преминаване от 16-байтови на
32-байтови вектори. Тя обаче не е двукратна, каквато би била при чисто аритметично
ограничение. Причината е, че вторият цикъл, този с проверките и обновяването на
купчината, не се векторизира и не се променя от ширината на регистъра.

Редът с лист от 1\,024 точки и $k = 8$ се отличава: отношението е 1,33 и 1,30 срещу над 2
навсякъде другаде. При малък лист и голямо $k$ купчината се пълни и остава пълна, тоест
делът на невекторизираната работа е най-голям точно там. Това е същият механизъм, който
обяснява и Табл.~\ref{tab:knn}.

\subsection{Заявка по радиус и размер на листа}

Табл.~\ref{tab:radius} дава заявката по радиус. Средният брой намерени съседи е
детерминиран и се повтаря точно между изпълненията.

\begin{table}[H]
\centering
\caption{Заявка по радиус, 2\,000 заявки, подразбираща се конфигурация}
\label{tab:radius}
\begin{tabular}{@{}rrrrrr@{}}
\toprule
\textbf{Точки} & \textbf{Радиус} & \textbf{Скал., ms} & \textbf{Пак., ms} & \textbf{Отн.} & \textbf{Средно съседи} \\
\midrule
10\,000    & 0{,}20 &  0{,}52 &  0{,}44 & 1{,}19 &   0{,}3 \\
10\,000    & 0{,}50 &  1{,}01 &  0{,}91 & 1{,}11 &   4{,}9 \\
50\,000    & 0{,}20 &  0{,}85 &  0{,}96 & 0{,}88 &   1{,}6 \\
50\,000    & 0{,}50 &  2{,}37 &  2{,}88 & 0{,}82 &  24{,}5 \\
200\,000   & 0{,}20 &  2{,}56 &  1{,}72 & 1{,}49 &   6{,}5 \\
200\,000   & 0{,}50 &  8{,}40 &  6{,}02 & 1{,}40 &  98{,}5 \\
1\,000\,000 & 0{,}20 &  6{,}25 &  6{,}05 & 1{,}03 &  32{,}6 \\
1\,000\,000 & 0{,}50 & 32{,}81 & 43{,}13 & 0{,}76 & 494{,}0 \\
\bottomrule
\end{tabular}
\end{table}

Отношението отново се движи около единица и не подкрепя извод. За заявката по радиус това
е и очакваното: при нея скъпата част не е сметката на разстоянието, а записването на
намерените съседи в изходния вектор, а то е еднакво и за двата пътя. При радиус 0,50 и
милион точки на заявка се връщат средно 494 съседа, тоест почти цялата работа е запис.

Табл.~\ref{tab:leaf} дава влиянието на размера на листа при 200\,000 точки и $k = 8$.

\begin{table}[H]
\centering
\caption{Размер на листа, 200\,000 точки, $k = 8$, 5\,000 заявки}
\label{tab:leaf}
\begin{tabular}{@{}rrrrrrrr@{}}
\toprule
& \multicolumn{3}{c}{\textbf{Подразбираща се}} & & \multicolumn{3}{c}{\textbf{\code{-march=native}}} \\
\cmidrule(lr){2-4}\cmidrule(lr){6-8}
\textbf{Лист} & \textbf{Постр., ms} & \textbf{Скал., ms} & \textbf{Пак., ms} & &
\textbf{Постр., ms} & \textbf{Скал., ms} & \textbf{Пак., ms} \\
\midrule
4   & 66{,}03 & 13{,}18 & 13{,}62 & & 71{,}52 & 16{,}71 & 14{,}69 \\
8   & 53{,}90 &  9{,}66 & 10{,}27 & & 60{,}67 & 14{,}05 & 10{,}71 \\
16  & 49{,}84 &  8{,}86 &  8{,}71 & & 48{,}26 &  8{,}80 &  8{,}63 \\
32  & 54{,}84 & 16{,}79 & 17{,}85 & & 46{,}51 &  9{,}59 & 12{,}35 \\
64  & 74{,}99 & 17{,}22 & 11{,}21 & & 72{,}41 & 12{,}05 & 10{,}67 \\
128 & 56{,}69 & 15{,}43 & 15{,}48 & & 42{,}37 & 14{,}24 & 12{,}20 \\
\bottomrule
\end{tabular}
\end{table}

Двете конфигурации се съгласяват за най-добрия размер на листа: 16 точки, с 8,71 ms и
8,63 ms за пакетния път, съответно 8,86 ms и 8,80 ms за скаларния. Съгласието на две
независими изпълнения върху една и съща стойност е по-силен довод от всяко отделно число
в таблицата. При лист 4 работата е в обхождането на дървото, при лист 128 в излишно
проверени точки, а между двете има минимум.

Времето за построяване спада с растежа на листа, което е очаквано: по-големият лист
означава по-плитко дърво и по-малко нива на разделяне.

\subsection{Прореждане с вокселна решетка}

Табл.~\ref{tab:voxel} дава прореждането на милион точки. Броят изходни точки е
детерминиран и съвпада точно между двете конфигурации и между изпълненията.

\begin{table}[H]
\centering
\caption{Прореждане с вокселна решетка, 1\,000\,000 точки в куб със страна 10}
\label{tab:voxel}
\begin{tabular}{@{}rrrrr@{}}
\toprule
\textbf{Ребро} & \textbf{Мин., ms} & \textbf{Мин., ms (native)} & \textbf{Изходни точки} & \textbf{Остават, \%} \\
\midrule
0{,}020 & 150{,}28 & 293{,}71 & 995\,925 & 99{,}59 \\
0{,}050 & 146{,}82 & 208{,}18 & 940\,245 & 94{,}02 \\
0{,}100 & 131{,}54 & 197{,}65 & 631\,816 & 63{,}18 \\
0{,}250 & 107{,}38 & 155{,}69 &  64\,000 &  6{,}40 \\
\bottomrule
\end{tabular}
\end{table}

Времето е доминирано от сортирането на милион ключа и слабо зависи от реброто: между
най-финото и най-едрото прореждане разликата е 1,4 пъти, докато броят изходни точки пада
15 пъти. Това съответства на устройството: сортирането се прави винаги върху целия вход, а
от реброто зависи само колко групи ще излязат от него.

Редът с ребро 0,250 е достоен за отбелязване: изходът е точно 64\,000 точки, тоест
$40^3$, и това е броят на кубчетата, а не на точките. При това ребро всяко кубче в куба
със страна 10 съдържа поне една от милиона точки, тоест решетката е наситена.

Числата в конфигурацията с \code{-march=native} са системно по-големи, което не се тълкува
като ефект от компилирането: прореждането няма векторизиран цикъл, а изпълнението с
\code{native} е било второто по ред и е попаднало на по-натоварена машина.

\subsection{Съвместяване}

Табл.~\ref{tab:icp} дава поведението на съвместяването спрямо нивото на шума. Всички
колони без последната са детерминирани и се повтарят точно между двете конфигурации.

\begin{table}[H]
\centering
\caption{Съвместяване, точка към точка. Праг за отхвърляне 0,5, най-много 60 итерации.
Зададеното преобразувание е ротация от 7,45$^\circ$ и транслация с дължина 0,229}
\label{tab:icp}
\begin{tabular}{@{}rrlrrr@{}}
\toprule
\textbf{Шум $\sigma$} & \textbf{Итерации} & \textbf{Сходимост} & \textbf{Отклонение} & \textbf{Грешка по ъгъл, $^\circ$} & \textbf{Мин., ms} \\
\midrule
0{,}000 & 38 & да & 0{,}015230 & 0{,}0055 & 309{,}01 \\
0{,}002 & 39 & да & 0{,}015365 & 0{,}0016 & 326{,}28 \\
0{,}005 & 41 & да & 0{,}017720 & 0{,}0035 & 380{,}62 \\
0{,}010 & 49 & да & 0{,}022771 & 0{,}0040 & 468{,}24 \\
0{,}020 & 60 & не & 0{,}030791 & 0{,}0430 & 641{,}42 \\
0{,}040 & 60 & не & 0{,}038900 & 0{,}8019 & 829{,}56 \\
\bottomrule
\end{tabular}
\end{table}

Броят итерации до сходимост расте с шума: 38, 39, 41, 49, и при $\sigma \ge 0{,}020$
границата от 60 итерации се изчерпва, преди условието за спиране да е изпълнено. Грешката
по ъгъл остава под 0,006$^\circ$ до $\sigma = 0{,}010$, тоест до шум, равен на една пета
от реброто на прореждащата решетка. При $\sigma = 0{,}040$ тя скача на 0,80$^\circ$, което
е двеста пъти повече, макар отклонението да е нараснало само 2,6 пъти. Двете величини
измерват различни неща и точно това разминаване го показва: отклонението мери колко близо
са точките, а грешката по ъгъл мери дали са в \term{правилното} положение, и второто се
разпада по-рязко.

Отклонението при $\sigma = 0$ не е нула, а 0,0152. Това не е грешка на процедурата, а
следствие от постановката: двата облака съдържат \term{различни} точки, извадка от 75\,\%
срещу пълния набор, прекарани през прореждане с ребро 0,05. Дори при съвършено съвместяване
най-близкият съсед на точка от единия облак е на разстояние от порядъка на разстоянието
между съседни точки в другия. Стойността 0,0152 е около една трета от реброто на решетката,
което е точно очакваният порядък.

Демонстрацията, чийто изход е независим от този експеримент, дава същото поведение при
$\sigma = 0{,}004$: 40 итерации, отклонение 0,016836, грешка по ъгъл 0,002035$^\circ$ и
грешка по транслация 0,000158 при транслация с дължина 0,229, тоест относителна грешка под
0,07\,\%.

\subsection{Растерен конвейер}

Табл.~\ref{tab:raster} дава времената на стъпките върху кадър 1024 на 768, тоест 786\,432
пиксела.

\begin{table}[H]
\centering
\caption{Растерен конвейер, кадър 1024 на 768, 5 повторения, подразбираща се конфигурация}
\label{tab:raster}
\begin{tabular}{@{}L{6.0cm}rrr@{}}
\toprule
\textbf{Стъпка} & \textbf{Мин., ms} & \textbf{Мед., ms} & \textbf{ns на пиксел} \\
\midrule
Гаусово изглаждане, $\sigma = 1{,}6$ & 47{,}305 & 52{,}532 & 60{,}2 \\
Градиент по Собел                    & 16{,}375 & 31{,}768 & 20{,}8 \\
Потискане на немаксимуми             & 11{,}225 & 11{,}739 & 14{,}3 \\
Двупрагово решение                   &  0{,}707 &  0{,}849 &  0{,}9 \\
Свързани компоненти                  &  5{,}925 &  5{,}960 &  7{,}5 \\
\bottomrule
\end{tabular}
\end{table}

Изглаждането доминира. При $\sigma = 1{,}6$ радиусът е 5, тоест ядрото е с 11 отчета и
разделимото изглаждане прави 22 умножения на пиксел. Градиентът по Собел прави четири
преминавания с ядро от 3 отчета, тоест 12 умножения на пиксел, плюс по един корен за
големината. Отношението на броя умножения е 1,8, а измереното отношение на времената е
47,305 към 16,375, тоест 2,9. Разликата между двете не се обяснява от броя операции и не се
обяснява тук: причината ѝ не е измерена, а предположение на нейно място би било число без
покритие.

Двупраговото решение е с два порядъка по-евтино от останалите стъпки, защото за повечето
пиксели се свежда до едно сравнение: проследяването тръгва само от пикселите над горния
праг, а те са малка част от кадъра. В демонстрацията, върху кадър 320 на 240, по контур
попадат 1\,223 пиксела, тоест 1,59\,\% от кадъра.

\subsection{Анализ}

Три извода са подкрепени от таблиците и един въпрос остава без отговор.

\textbf{Първо, векторизацията на обхождането на лист работи и печалбата ѝ е около
двукратна.} Табл.~\ref{tab:leafscan} я показва в дванадесет клетки при две независими
конфигурации на компилиране, с една и съща посока навсякъде. Че цикълът наистина се
превежда във векторни инструкции, е потвърдено и независимо от времето, чрез отчета на
компилатора, цитиран в Раздел~\ref{sec:implementation}.

\textbf{Второ, тази печалба не се вижда в цялата заявка.} Табл.~\ref{tab:knn} не позволява
извод: отношението се движи около единица, а разсейването надхвърля разликата, която се
търси. Обяснението следва от първите два извода взети заедно. Обхождането на дървото,
работата с купчината и записът на резултата не се векторизират, и колкото по-голямо е $k$,
толкова по-голям е делът им. Редът с най-малък лист и най-голямо $k$ в
Табл.~\ref{tab:leafscan}, където отношението пада от над 2 на 1,3, е същият механизъм,
видян отвътре.

Това е отклонение от очакваното в Раздел~\ref{sec:alternatives}, където разположението по
компонента беше избрано именно заради заявката. Отклонението се отчита като отклонение.
Изборът на разположение не е погрешен, той дава двукратно ускорение на частта, която
покрива, но частта е по-малка от предполаганото и затова печалбата в целия конвейер е
по-малка от очакваната.

\textbf{Трето, съвместяването възстановява зададеното преобразувание с грешка под
0,006$^\circ$ при шум до една пета от реброто на прореждащата решетка,} и се разпада
предвидимо над това ниво (Табл.~\ref{tab:icp}).

Въпросът, който остава без отговор, е количествен: каква точно е печалбата от
векторизацията в цялата заявка. Тя не е нула, защото частичната печалба е измерена, но е
под разделителната способност на машината, на която са проведени измерванията.
Отговорът изисква машина без чужд товар по ядрата, а не друга реализация. Това е
ограничение на експеримента и е записано като такова, вместо да бъде заместено с
правдоподобно число.
